\documentclass{beamer}
\mode<presentation>
\usepackage{amsmath}
\usepackage{amssymb}
\usepackage[utf8]{inputenc}
\usepackage{fontenc}
%\usepackage{advdate}
\usepackage{adjustbox}
\usepackage{subcaption}
\usepackage{enumitem}
\usepackage{multicol}
\usepackage{mathtools}
\usepackage{listings}
\usepackage{url}
\def\UrlBreaks{\do\/\do-}
\usetheme{Boadilla}
\usecolortheme{lily}
\setbeamertemplate{footline}
{
  \leavevmode%
  \hbox{%
  \begin{beamercolorbox}[wd=\paperwidth,ht=2.25ex,dp=1ex,right]{author in head/foot}%
    \insertframenumber{} / \inserttotalframenumber\hspace*{2ex} 
  \end{beamercolorbox}}%
  \vskip0pt%
}
\setbeamertemplate{navigation symbols}{}

\providecommand{\nCr}[2]{\,^{#1}C_{#2}} % nCr
\providecommand{\nPr}[2]{\,^{#1}P_{#2}} % nPr
\providecommand{\mbf}{\mathbf}
\providecommand{\pr}[1]{\ensuremath{\Pr\left(#1\right)}}
\providecommand{\qfunc}[1]{\ensuremath{Q\left(#1\right)}}
\providecommand{\sbrak}[1]{\ensuremath{{}\left[#1\right]}}
\providecommand{\lsbrak}[1]{\ensuremath{{}\left[#1\right.}}
\providecommand{\rsbrak}[1]{\ensuremath{{}\left.#1\right]}}
\providecommand{\brak}[1]{\ensuremath{\left(#1\right)}}
\providecommand{\lbrak}[1]{\ensuremath{\left(#1\right.}}
\providecommand{\rbrak}[1]{\ensuremath{\left.#1\right)}}
\providecommand{\cbrak}[1]{\ensuremath{\left\{#1\right\}}}
\providecommand{\lcbrak}[1]{\ensuremath{\left\{#1\right.}}
\providecommand{\rcbrak}[1]{\ensuremath{\left.#1\right\}}}
\theoremstyle{remark}
\newtheorem{rem}{Remark}
\newcommand{\sgn}{\mathop{\mathrm{sgn}}}
%\providecommand{\abs}[1]{\left\vert#1\right\vert}
\providecommand{\res}[1]{\Res\displaylimits_{#1}} 
\providecommand{\norm}[1]{\lVert#1\rVert}
\providecommand{\mtx}[1]{\mathbf{#1}}
%\providecommand{\mean}[1]{E\left[ #1 \right]}
\providecommand{\fourier}{\overset{\mathcal{F}}{ \rightleftharpoons}}
%\providecommand{\hilbert}{\overset{\mathcal{H}}{ \rightleftharpoons}}
\providecommand{\system}{\overset{\mathcal{H}}{ \longleftrightarrow}}
	%\newcommand{\solution}[2]{\textbf{Solution:}{#1}}
%\newcommand{\solution}{\noindent \textbf{Solution: }}
\providecommand{\dec}[2]{\ensuremath{\overset{#1}{\underset{#2}{\gtrless}}}}
\newcommand{\myvec}[1]{\ensuremath{\begin{pmatrix}#1\end{pmatrix}}}
\let\vec\mathbf

\lstset{
language=C,
frame=single, 
breaklines=true,
columns=fullflexible
}

\numberwithin{equation}{section}

\title{4.2.1}
\author{AI25BTECH11038 -- Tejas Uppala}

\begin{document}

\maketitle
\begin{frame}{Question}
  Find the direction and normal vectors of the following line $2x + 3y = 9$  
\end{frame}

\begin{frame}{Solution}
  For a general line y = mx + c, it can be converted into its vector form in the following way,  
  \begin{equation}
    y = mx + c
\end{equation}

\begin{equation}
    \myvec{x\\y} = \myvec{x\\mx+c} = \myvec{0\\c} + x\myvec{1\\m}    
\end{equation}
yielding, 
\begin{equation}
    \textbf{x} = \textbf{h} + k\textbf{m}
\end{equation}
\end{frame}

\begin{frame}
    where \textbf{h} is any point on the line and \textbf{m} is the direction vector of the line,

\begin{equation}
    \textbf{m} = \myvec{1\\m}
\end{equation}

The given line can be written as, 

\begin{equation}
    y = \brak{\frac{-2}{3}}x + \frac{9}{3} = \brak{\frac{-2}{3}}x + 3
\end{equation}

on correlating,

\begin{equation}
    \textbf{m} = \myvec{1 \\ \frac{-2}{3}} or \myvec{3\\-2}, \textbf{h} = \myvec{0\\3}
\end{equation}

The \textbf{normal vector} (n) of the line is \myvec{-m\\1}. Therefore, 

\begin{equation}
    \textbf{n} = \myvec{\frac{2}{3}\\1} or \myvec{2\\3}
\end{equation}
\end{frame}


\end{document}

\end{document}
